% RL in Production — Capstone Project: The RLVR Arena
% Vizuara AI Labs · compile with XeLaTeX
%   xelatex capstone-rlvr-arena.tex   (run twice)
\documentclass[11pt]{article}
\usepackage{fontspec}
\usepackage[a4paper, margin=1.9cm, top=2.0cm, bottom=2.2cm]{geometry}
\usepackage{xcolor}
\usepackage{amsmath}
\usepackage{amssymb}
\usepackage[most]{tcolorbox}
\usepackage{enumitem}
\usepackage{titlesec}
\usepackage{microtype}
\usepackage{fancyvrb}
\usepackage{graphicx}
\usepackage{booktabs}
\usepackage{array}
\usepackage{fancyhdr}
\usepackage{hyperref}

% ── Palette ──────────────────────────────────────────────────────────────────
\definecolor{cream}{HTML}{F7F3EC}
\definecolor{paper}{HTML}{FBF8F2}
\definecolor{ink}{HTML}{2B2B2B}
\definecolor{ink2}{HTML}{55514A}
\definecolor{ink3}{HTML}{8A8475}
\definecolor{indigo}{HTML}{4C5B9A}
\definecolor{amber}{HTML}{C8851B}
\definecolor{mint}{HTML}{4E8D6E}
\definecolor{rose}{HTML}{B5566B}
\definecolor{plum}{HTML}{6B4E7E}
\pagecolor{cream}\color{ink}

% ── Fonts ────────────────────────────────────────────────────────────────────
\setmainfont{Charter}[Numbers=OldStyle, Ligatures=TeX]
\newfontfamily\sans{Avenir Next}[Ligatures=TeX]
\newfontfamily\sansbold{Avenir Next}[UprightFont={* Bold}, Ligatures=TeX]
\newfontfamily\sanstrack{Avenir Next}[LetterSpace=14.0, Ligatures=TeX]
\setmonofont{Menlo}[Scale=0.78]

% ── Commands ─────────────────────────────────────────────────────────────────
\newcommand{\eyebrow}[1]{{\sanstrack\footnotesize\color{ink3}\MakeUppercase{#1}}\par\vspace{0.2em}}
\newcommand{\pts}[1]{{\sans\footnotesize\color{amber}\textbf{[#1]}}}
\newcommand{\provided}{{\sans\footnotesize\color{mint}\textbf{[\,PROVIDED\,]}}}
\newcommand{\todotag}[1]{{\sans\footnotesize\color{rose}\textbf{[\,TODO #1\,]}}}
\newcommand{\lbtag}{{\sans\footnotesize\color{plum}\textbf{[\,LEADERBOARD\,]}}}
\newcommand{\codesm}[1]{{\ttfamily\footnotesize #1}}

% ── Section style ─────────────────────────────────────────────────────────────
\titleformat{\section}{\sansbold\Large\color{ink}}{}{0pt}{}[\vspace{-0.4em}{\color{ink3}\rule{\textwidth}{0.4pt}}]
\titleformat{\subsection}{\sansbold\normalsize\color{ink2}}{}{0pt}{}
\titlespacing{\section}{0pt}{1.4em}{0.5em}
\titlespacing{\subsection}{0pt}{0.9em}{0.3em}
\setlength{\parindent}{0pt}\setlength{\parskip}{0.48em}\linespread{1.07}

% ── Boxes ────────────────────────────────────────────────────────────────────
\newtcolorbox{qbox}[1]{enhanced, colback=paper, colframe=#1, boxrule=0pt,
  leftrule=3pt, arc=2mm, left=5mm, right=5mm, top=2.5mm, bottom=2.5mm, breakable}
\newtcolorbox{deliver}{enhanced, colback=mint!8, colframe=mint, boxrule=0pt,
  leftrule=3pt, arc=2mm, left=5mm, right=5mm, top=2.2mm, bottom=2.2mm, breakable}
\newtcolorbox{warn}{enhanced, colback=amber!7, colframe=amber, boxrule=0pt,
  leftrule=3pt, arc=2mm, left=5mm, right=5mm, top=2.2mm, bottom=2.2mm, breakable}
\newtcolorbox{herobox}{enhanced, colback=indigo!7, colframe=indigo, boxrule=0pt,
  leftrule=3pt, arc=2mm, left=6mm, right=6mm, top=3mm, bottom=3mm, breakable}
\DefineVerbatimEnvironment{code}{Verbatim}{fontsize=\footnotesize, frame=leftline,
  framerule=0.6pt, rulecolor=\color{ink3}, framesep=7pt, baselinestretch=0.96}

% ── Header / footer ──────────────────────────────────────────────────────────
\pagestyle{fancy}\fancyhf{}
\renewcommand{\headrulewidth}{0.3pt}
\fancyhead[L]{\sans\footnotesize\color{ink3} RL in Production · Capstone Project}
\fancyhead[R]{\sans\footnotesize\color{ink3} Vizuara AI Labs}
\fancyfoot[C]{\sans\footnotesize\color{ink3}\thepage}

\hypersetup{colorlinks=true,linkcolor=indigo,urlcolor=indigo,
  pdftitle={RL Capstone — The RLVR Arena},pdfauthor={Vizuara AI Labs}}

% ─────────────────────────────────────────────────────────────────────────────
\begin{document}
\thispagestyle{empty}

\vspace*{0.3cm}
\eyebrow{RL in Production \ \textbullet\ \ Vizuara AI Labs \ \textbullet\ \ Capstone Project}
\vspace{0.3cm}

{\fontsize{28}{33}\selectfont \textbf{The RLVR Arena}\par}
\vspace{0.2cm}
{\fontsize{14}{18}\selectfont\color{ink2}\itshape
Train a language model with GRPO to solve Countdown number puzzles.
Submit your checkpoint. It gets re-run on a secret test set.
The result goes on a leaderboard.\par}

\vspace{0.5cm}{\color{ink3}\rule{\textwidth}{0.4pt}}\vspace{0.3cm}

\begin{herobox}
\textbf{The throughline of this entire series closes here.}
Every lecture traced one recursion: $V(s) = r + \gamma V(s')$.
We solved it with dynamic programming (L1), approximated it with a neural network (L2),
climbed it directly with policy gradients (L3), stabilised the climb with PPO's clip (L4),
and applied it to a language model with RLHF (L5).
In L6 we stripped away the critic entirely and replaced it with a group of sampled answers ---
that was GRPO.
Now you implement it for real, on a real model, on a real verifiable task,
and compete on a leaderboard.
\end{herobox}

% ─────────────────────────────────────────────────────────────────────────────
\section{The Task --- Countdown}

You are given a list of numbers and a target.
Combine \textbf{all} the numbers --- each used exactly once --- with $+\,-\,\times\,\div$ and parentheses to reach the target.

\begin{qbox}{indigo}
\textbf{Example.} \quad Numbers: \textbf{3, 7, 8, 2} \quad Target: \textbf{24}

\vspace{0.3em}
\hspace{2em}$\underbrace{(8 - 2)}_6 \times \underbrace{(7 - 3)}_4 = 24$ \quad\checkmark
\quad uses 8, 2, 7, 3 exactly once.

\vspace{0.3em}
Your model must reason inside \codesm{<think>\,\ldots\,</think>} and give its final expression inside \codesm{<answer>\,\ldots\,</answer>}.
That is the entire output contract.
\end{qbox}

\textbf{Why this task?}
Correctness is a five-line arithmetic checker --- no LLM judge, no rubric, no human in the loop.
The task is combinatorial so there is nothing to memorise from the training set.
And a 0.5B model that has never seen this task goes from near-zero to genuine reasoning under GRPO ---
that is the ``aha moment'' you are here to produce.

% ─────────────────────────────────────────────────────────────────────────────
\section{What You Implement}

Open \codesm{train\_grpo.py}. Two functions are marked \todotag{}: everything else is \provided.

\subsection{TODO 1 --- Group-Relative Advantage \pts{12}}

\begin{qbox}{rose}
\textbf{Input.} A flat tensor of rewards of shape $[\textit{num\_puzzles} \times \textit{group}]$,
where the $k$ completions for each puzzle are grouped consecutively.

\textbf{Output.} An advantage tensor of the same shape.

The single idea of L6: within each group, how much better (or worse) was this answer compared to
the \emph{rest of the group}?
That relative score --- not the absolute reward --- is the advantage.
No critic. No baseline model. Just the group.

Review the L6 slides before writing this function.
\end{qbox}

\subsection{TODO 2 --- GRPO Loss \pts{13}}

\begin{qbox}{rose}
\textbf{Inputs.} Per-token log-probs under the policy ($\log\pi_\theta$),
the behaviour policy ($\log\pi_{\text{old}}$), and the frozen reference ($\log\pi_{\text{ref}}$);
a completion mask; and the sequence-level advantages from TODO~1.

\textbf{Output.} A scalar loss.

Two ideas from lecture combine here:
\begin{itemize}[itemsep=1pt]
  \item \textbf{L4 (PPO):} the clipped importance-sampling surrogate.
        Compute per-token probability ratios and clip them so no single step moves too far.
  \item \textbf{L5 (RLHF):} a KL penalty to the frozen reference keeps the model from
        drifting away from its starting point.
\end{itemize}
Average the per-token objective over actual completion tokens (use the mask).
Return the \emph{negative} (you minimise; gradient ascent happens automatically).
The KL estimator to use is the k3 form from the GRPO paper:
$\exp(\log\pi_{\text{ref}} - \log\pi_\theta) - (\log\pi_{\text{ref}} - \log\pi_\theta) - 1$.
\end{qbox}

\textbf{That is all.}
The rollout (sampling a group of completions and computing rewards), the per-token log-prob helper,
the training loop, checkpointing, and evaluation harness are fully provided.

% ─────────────────────────────────────────────────────────────────────────────
\section{Files You Are Given}

\begin{center}
\renewcommand{\arraystretch}{1.35}
\begin{tabular}{>{\ttfamily\footnotesize}l p{10.2cm}}
\toprule
\normalfont\textbf{File} & \textbf{Purpose} \\
\midrule
countdown.py     & The task engine. Contains the exact verifier used to score your submission.
                   Read it --- especially \codesm{verify()} and \codesm{reward()}. \\
train\_grpo.py   & GRPO training starter. Fill the two \todotag{} functions. Everything else runs. \\
run\_model.py    & Turn a saved checkpoint into \codesm{predictions.jsonl} for self-scoring. \\
evaluate.py      & Score a predictions file $\to$ accuracy breakdown. Run this on \codesm{dev\_public}
                   to see your number before submitting. \\
data/train.jsonl & 4\,000 puzzles to train on (easy / medium / hard mix). \\
data/dev\_public.jsonl & 300 puzzles to self-score.
                   The leaderboard uses the \emph{same metric} on a larger held-out set. \\
requirements.txt & Python dependencies. \\
\bottomrule
\end{tabular}
\end{center}

\begin{qbox}{mint}
\textbf{Quick start.}
\begin{code}
pip install -r requirements.txt

# Smoke-test (CPU, tiny model, 2 steps -- proves your code runs):
python train_grpo.py --model sshleifer/tiny-gpt2 \
    --steps 2 --group 4 --bsz 2 --smoke

# Self-score on the public dev set:
python run_model.py --model ./model \
    --test data/dev_public.jsonl --out dev_preds.jsonl
python evaluate.py --test data/dev_public.jsonl \
    --predictions dev_preds.jsonl --name yourname
\end{code}
\end{qbox}

% ─────────────────────────────────────────────────────────────────────────────
\section{The Base Model and the Rules}

\begin{center}
\renewcommand{\arraystretch}{1.32}
\begin{tabular}{l p{10.5cm}}
\toprule
\textbf{Rule} & \textbf{Detail} \\
\midrule
Base model & \codesm{Qwen/Qwen2.5-0.5B-Instruct} (frozen starting point, no exceptions) \\
Method & GRPO or any policy-gradient RL you implement. \textbf{No SFT on solution traces}
         --- the dataset contains only puzzles, no worked answers. \\
No tools & The model reasons in tokens only. No calculator, no code execution, no API calls. \\
Eval config & Greedy decode, \codesm{max\_new\_tokens=640}, the canonical prompt in \codesm{countdown.py} ---
             identical for everyone. Deterministic $\Rightarrow$ reproducible $\Rightarrow$ fair. \\
Open track & Optional: any base model $\leq 3$B, ranked on a separate board.
             Your reward curve must show GRPO is doing the work. \\
\bottomrule
\end{tabular}
\end{center}

\begin{warn}
\textbf{Fixed base model = fixed starting line.}
The leaderboard measures your \emph{RL}, not who has the bigger GPU or a stronger base.
Everyone starts from exactly the same 0.5B checkpoint.
\end{warn}

% ─────────────────────────────────────────────────────────────────────────────
\section{The Leaderboard \lbtag}

You submit a checkpoint. We load it, run \codesm{run\_model.py} on a \textbf{private test set you have never seen}, and \codesm{evaluate.py} produces the official score. Your self-reported accuracy is never trusted --- only the re-run counts.

\begin{center}
\renewcommand{\arraystretch}{1.32}
\begin{tabular}{l l p{7.5cm}}
\toprule
\textbf{Rank key} & \textbf{Metric} & \textbf{Notes} \\
\midrule
\textbf{Primary} \lbtag & Accuracy on full private test & The headline number \\
Tie-break 1 & Accuracy on \textit{hard} split (5 numbers) & Rewards genuine reasoning \\
Tie-break 2 & Avg tokens when correct (fewer wins) & Efficiency, no rambling \\
Diagnostic & Format compliance rate & Did it learn the protocol? \\
\bottomrule
\end{tabular}
\end{center}

\textbf{Reference points seeded on the board:} (1) the untrained base model (the floor --- beat this to pass), and (2) a reference GRPO run (the target --- beat this to compete).

\begin{qbox}{plum}
\textbf{The cold-start warning.}
GRPO typically shows \emph{near-zero accuracy for the first 400--600 steps} --- this is normal.
The model is learning the output format before it can reason.
The ``aha moment'' (reward suddenly climbing) arrives between step 600 and 1\,000 for most correct implementations.
Do not panic if your curve is flat early.
Run longer if you can.
\end{qbox}

% ─────────────────────────────────────────────────────────────────────────────
\section{What to Submit}

\begin{deliver}
\begin{enumerate}[leftmargin=1.5em, itemsep=3pt]
  \item \textbf{\codesm{model/}} --- your fine-tuned checkpoint.
        Must load with \codesm{AutoModelForCausalLM.from\_pretrained("./model")}.
        Full weights \emph{or} a LoRA adapter with the base model ID.
  \item \textbf{\codesm{dev\_preds.jsonl}} --- your predictions on \codesm{dev\_public.jsonl},
        generated by \codesm{run\_model.py}. Used for a reproducibility check:
        re-running your checkpoint must give the same accuracy within~$\pm 0.3\%$.
  \item \textbf{\codesm{train\_grpo.py}} --- your filled-in training code.
        The TODOs must be implemented; the rest may be modified.
  \item \textbf{\codesm{reward\_curve.png}} --- a plot of mean reward vs.\ training step.
        Export it from your training log.
  \item \textbf{Report ($\leq 1$ page)} --- what you tried, the reward curve interpreted,
        one ablation (e.g.\ group size, KL $\beta$, learning rate), one failure mode you encountered.
\end{enumerate}
\end{deliver}

% ─────────────────────────────────────────────────────────────────────────────
\section{Grading}

\begin{center}
\renewcommand{\arraystretch}{1.35}
\begin{tabular}{l r}
\toprule
\textbf{Component} & \textbf{Points} \\
\midrule
Working checkpoint: loads, runs, and beats the untrained floor & 40 \\
GRPO implemented correctly (both TODOs, read in your code) & 25 \\
Report: reward curve + one ablation + one failure mode & 20 \\
Reproducibility: \codesm{dev\_preds} match our re-run of your checkpoint & 10 \\
Leaderboard placement (top 3 gets full 5; everyone who submits gets 3) & 5 \\
\midrule
\textbf{Total} & \textbf{100} \\
\bottomrule
\end{tabular}
\end{center}

% ─────────────────────────────────────────────────────────────────────────────
\section{Compute}

\begin{qbox}{indigo}
\textbf{Colab T4 (free).}
The 0.5B model fits on a 16\,GB T4 at \codesm{float16}.
Use \codesm{--group 4 --bsz 4} (16 sequences per step).
Call \codesm{model.gradient\_checkpointing\_enable()} after loading to halve activation memory.
Checkpoint to Drive to survive disconnects.
One full run at 1\,000 steps takes roughly 2--3 hours.

\vspace{0.4em}
\textbf{Modal (if you have an account).}
Create a free Modal account at \codesm{modal.com}, run \codesm{modal token new} once, then use
\codesm{A10G} (\codesm{gpu="A10G"}) with \codesm{--group 8 --bsz 2}.
A 1\,000-step run finishes in about 2 hours and costs roughly \$1--2 in GPU credits.
\end{qbox}

% ─────────────────────────────────────────────────────────────────────────────
\section{Timeline}

\begin{center}
\renewcommand{\arraystretch}{1.32}
\begin{tabular}{l p{11cm}}
\toprule
\textbf{Day} & \textbf{Milestone} \\
\midrule
1 & Run the smoke-test (CPU, tiny-gpt2, 2 steps). Read \codesm{countdown.py} end to end.
    Run the \emph{untrained} base model through \codesm{evaluate.py} on dev --- see the floor. \\
2--3 & Implement both TODOs. Run on easy puzzles only (\codesm{data/train.jsonl} filtered).
       Get the reward moving. \\
4--5 & Full difficulty mix. Tune: group size, KL $\beta$, learning rate. Log everything. \\
6 & Lock a checkpoint. Generate \codesm{dev\_preds.jsonl}. Write the report. \\
7 & Submit. Leaderboard goes live. \\
\bottomrule
\end{tabular}
\end{center}

% ─────────────────────────────────────────────────────────────────────────────
\section*{One last thing}

\vfill
The verifier in \codesm{countdown.py} uses \emph{exact rational arithmetic} (\codesm{fractions.Fraction})
and a whitelisted-AST evaluator --- it never calls \codesm{eval()} on your model's output.
A model that tries to write \codesm{\_\_import\_\_('os').system('...')} inside its \codesm{<answer>}
tag gets scored \emph{invalid}, not executed.
This is a good design principle to carry forward.

Good luck.
The throughline was $V(s) = r + \gamma V(s')$.
Now make it real.

\vspace{0.5cm}
{\color{ink3}\rule{\textwidth}{0.4pt}}
\vspace{0.1cm}
{\sans\footnotesize\color{ink3}
All code lives at \texttt{github.com/VizuaraAI/RL-in-Production-Bootcamp-Resources/lectures/capstone/}.
Questions in the course channel. \quad Vizuara AI Labs · 2026.}

\end{document}
